\chapter{Chapter 1 --- Introduction}\label{chapter-1-introduction}

This thesis studies how to \textbf{accelerate test-time discovery in complex code generation via execution-guided self-distillation}. The core question is practical: when a language model is given a single hard programming problem and a budget to improve itself at inference time, what is the most effective way to turn execution feedback into a better solution, and when does that strategy work?

\section{1.1 Background}\label{background}

Reinforcement learning with verifiable rewards (RLVR), with GRPO {[}8{]} as a common instantiation, suffers from sparse credit assignment: a scalar outcome reward says only whether an attempt passed, and when every rollout fails the advantage collapses and learning stalls --- exactly on the hard problems that matter. Self-Distillation Policy Optimization (SDPO) {[}1{]} extracts a denser signal from the rich feedback verifiable environments already produce, using the feedback-conditioned model as a self-teacher distilled back into the student. At test time it can iterate on one hard problem and improve before the first success {[}1, §5{]}. Two issues motivate this work: the parent paper distills the student's own failed rollout (little that is correct to learn from on hard problems), and Kim et al.~{[}2{]} show that distilling from rich context can degrade a model by suppressing its epistemic verbalization.

\section{1.2 Problem}\label{problem}

The setting is test-time training: one hard problem, a verifiable reward, a short budget of self-improvement steps, then reset. Two design choices govern discovery: \emph{how the execution feedback is organized into a learning signal} (learn from the student's failed attempt, or from a correct attempt produced under feedback and filtered for independence), and \emph{how the execution feedback is formulated} into the reprompt the self-teacher sees {[}1, §7{]}. We take execution feedback as the privileged context, hence \emph{execution-guided}, and study both choices with the goal of accelerating discovery.

\section{1.3 Research questions}\label{research-questions}

\begin{itemize}
\tightlist
\item
  \textbf{RQ-method.} Does a teacher-first organization beat the student-first baseline at discovering solutions to hard code problems, and does the advantage survive at matched compute?
\item
  \textbf{RQ1.} Does the reprompt-template formulation of the execution feedback affect discovery, and in which regime?
\item
  \textbf{RQ2.} Does self-distillation from rich context suppress epistemic verbalization at test time on code?
\item
  \textbf{RQ3.} How does discovery trade off against compute (compute-to-correct)?
\end{itemize}

All four research questions are systematically evaluated and addressed throughout the empirical framework of this thesis.

\section{1.4 Contributions}\label{contributions}

\begin{enumerate}
\def\labelenumi{\arabic{enumi}.}
\tightlist
\item
  \textbf{A teacher-first variant of execution-guided self-distillation} that distills the student toward verifier- and judge-filtered teacher generations, between on-policy SDPO and off-policy SFT, with a clean filter that never distils a copy (Chapter 3).
\item
  \textbf{Teacher-first significantly accelerates discovery, and the advantage is compute-fair.} Across a medium matched evaluation, teacher-first beats student-first (Wilcoxon signed-rank) and escapes the flat-reward trap on the hardest problems; a compute-matched comparison and a compute-to-correct frontier show the gain is not a sampling-volume artifact (Chapter 4).
\item
  \textbf{A reprompt-template effect}, strongest on hard problems, where a reasoning-inducing template orders above the anchor, which orders above a minimal one (§4.3).
\item
  \textbf{A measured epistemic-suppression result on code} (thinking enabled): distillation from answer-bearing context suppresses uncertainty markers, accumulating across steps (§4.5).
\item
  \textbf{A validated value-versus-procedure boundary.} Supplying a method-bearing reference (MATH-500 worked solutions) lets the model escape on math where an answer-only reference does not, confirming that escape requires an in-reach method, not merely a value (§4.6, Chapter 5).
\end{enumerate}

\section{1.5 Scope}\label{scope}

Test-time training with LoRA adapters; code generation on LiveCodeBench v6 {[}6{]} as the core, with AIME 2026 and MATH-500 as the math probes. Results are on one 4B-scale model. The study remains an empirical characterization in one model family at a personal compute scale, establishing foundational insights rather than universal scaling laws.

\section{1.6 Thesis structure}\label{thesis-structure}

Chapter 2 reviews background and related work; Chapter 3 specifies the method (including the fallback-free filter and the thinking-on and method-reference extensions); Chapter 4 reports the medium-scale experiments, the compute-matched comparison, the epistemic-suppression result, and the validated boundary; Chapter 5 discusses the mechanism; Chapter 6 states the remaining limitations and the publication roadmap; Chapter 7 concludes.